This section situates the present work with respect to broad strands of research that motivate graph-aware, citation-conscious generation for long-form scientific text. Because a comprehensive literature retrieval was not available at the time of drafting, the discussion below is intentionally high-level and framed as general background knowledge; a complete, citation-rich survey will be provided in the final manuscript once curated references are obtained.\footnote{Source: Plan: Literature (evidence pack); see note in manuscript metadata indicating that no online related work retrieval was performed for this draft.} The aim here is to outline the conceptual space and clarify how our design choices are informed by recurring themes in prior work, without asserting detailed provenance for specific ideas. We therefore focus on broad methodological categories and on describing motivations and desiderata that guided our design, rather than on cataloguing prior systems.

First, research on controllable text generation has explored mechanisms for steering model outputs via structured inputs, auxiliary control codes, or constrained decoding. These mechanisms are commonly used to influence high-level attributes such as style, length, and topical focus, and they are relevant to any system that must produce section-level scientific prose. In practice, such control techniques are applied both at generation time (e.g., via prompts or control tokens) and via architectural interventions (e.g., conditioning layers or adapters) to achieve predictable behaviors. The principal idea that motivates our use of schemas and scheduler constraints is the same: make intended structure explicit so that generation aligns with external requirements. The statement above is presented here as general background knowledge rather than a citation-backed literature claim, and we refrain from enumerating specific prior systems until a curated bibliography is assembled.

Second, retrieval-augmented generation (RAG) approaches combine an external evidence store with a neural generator so that produced claims can be grounded in retrieved documents or passages. The workflow presented in this paper adopts the same broad principle—injecting retrieved evidence into section-level prompts—but differs in emphasizing an explicit document graph representation and strict, schema-constrained section outputs. Concretely, our pipeline treats retrieval as one component in a larger orchestration that includes graph-based dependency tracking and output validation, rather than as a simple augmenting context. We also place operational emphasis on tracking which retrieved passages were used by which section and how those passages are connected across sections via the graph. This comparison is provided as general background context; detailed, citation-supported comparisons to canonical RAG instantiations will appear in the evaluation and final related-work revision.

Third, there has been growing interest in structured prompting and schema-constrained outputs, where models are required to emit machine-parseable artifacts (for example, JSON or XML) to facilitate downstream verification and assembly. Our use of strict JSON schemas for section outputs and the inclusion of runtime validation checks follows this practical desideratum and is intended to improve auditability at section and claim granularity. By enforcing a predictable output shape we aim to make automated downstream processing—such as extracting claims, mapping citation slots, and assembling a LaTeX document—more robust to model variability. In addition, schema constraints make it possible to run lightweight, automated consistency checks (for example, verifying that every citation slot is populated or that provenance links are non-empty) before accepting generated content. This paragraph is likewise framed as general background knowledge in the absence of a curated bibliography in the current draft.

Fourth, hierarchical and planner-based document generation techniques—ranging from top-down outline expansion to mixed bottom-up/top-down schedulers—aim to preserve coherence across lengthy documents by modeling dependencies among sections. The document-graph representation and the scheduler proposed in this paper are motivated by these design goals; they explicitly encode inter-section edges and temporal ordering constraints so that local generation decisions take global structure into account. We also emphasize mechanisms for propagating provenance and constraints along graph edges to reduce contradiction and manage cross-references. In practice this means that when a section cites a claim or a passage, that linkage is recorded on the connecting edge and can influence later generation steps; similarly, revision of an upstream node can trigger revalidation of dependent nodes. Precise comparisons to extant planner architectures will be provided once a targeted literature retrieval and citation curation are completed.

Fifth, systems explicitly designed for citation-aware generation seek to increase citation coverage and reduce hallucination by enforcing or encouraging links between assertions and source passages, and often produce provenance metadata alongside generated text. The present work shares this objective but places additional emphasis on an audit trail embedded in section-level JSON outputs and on an assembly mechanism that maps section citation slots to a refs.bib file during LaTeX generation. In addition to linking claims to passages, our approach records the mapping between in-text citation identifiers and bibliography entries to facilitate downstream reproducibility checks and manual inspection. We also record simple validation diagnostics (e.g., unreferenced bibliography entries, empty citation slots) in the JSON output so they can be surfaced to a human reviewer. This methodological distinction is descriptive of our design priorities rather than a claim of precedence.

Limitations of this section: no online related-work retrieval was performed for this draft, and therefore the authors are currently unable to confirm novelty or to fully position the contribution relative to specific prior publications.\footnote{Positioning note from evidence pack: ``Unable to confirm novelty due to missing related work.''; novelty risk: ``No online related work retrieved.''} We acknowledge that a careful, citation-rich comparison is necessary to validate the novelty and to surface close antecedents; such a comparison will include targeted citations, critical analysis of differences in objectives and evaluation methods, and an explicit mapping of our components to similar modules in prior systems. Concretely, the revised related work will identify representative systems in controllable generation, RAG, structured prompting, hierarchical planners, and citation-aware pipelines, and will discuss similarities and divergences in objectives, interface choices, and evaluation protocols. A complete literature review with explicit citations, critical comparison to closely related systems, and discussion of antecedent graph- or hierarchy-based document planners will be incorporated in the revised submission once curated references are collected.